\chapter*{Abstract}
\addcontentsline{toc}{chapter}{Abstract}

\vspace*{2cm}


This project focuses on the optimization and compression of Small Language Models
for Edge AI applications. The main objective is to prepare and fine-tune a compact
instruction-following language model, \texttt{google/flan-t5-small}, and to study
the impact of different optimization algorithms on training quality and resource
efficiency.

\vspace*{0.15cm}

The dataset used in this work is \texttt{yahma/alpaca-cleaned}, which was processed
into training, validation, and test subsets. Each example was transformed into an
instruction-tuning format composed of a source prompt and a target answer. The
fine-tuning pipeline was then designed to train the model using several optimizers,
including Stochastic Gradient Descent, SGD with momentum, Adam, and L-BFGS. Newton-CG
is also discussed from a theoretical perspective because it is not practical for
standard transformer fine-tuning in this context.

\vspace*{0.15cm}

The comparison is based on criteria that are important for Edge AI deployment:
validation loss, perplexity, training time, training speed, memory consumption,
model size, inference time, latency, and energy consumption when power telemetry is
available. These metrics make it possible to analyze the tradeoff between model
quality and computational efficiency.

\vspace*{0.15cm}

This report presents the project context, the data preparation process, the
fine-tuning methodology, the optimization algorithms studied, and the evaluation
strategy. The final goal is to identify which optimization approach provides the
best compromise for adapting a Small Language Model under the constraints of Edge
AI environments.
